%=============================================================================
% Wave 4, Iteration 10: Cosmology Deep Dive
%
% Agent: Cosmology Agent (W4i10, Opus 4.6)
% Date: 20 March 2026
%
% PRIORITIES:
%   1. T4 Cold Spot: Environmental MOND screening, full void ISW (50%)
%   2. alpha^57 Step 9: Coleman--Weinberg on CP2 x S3 (30%)
%   3. DESI DR2: Observer-ready predictions (20%)
%
% BUILDING ON:
%   W3i9_T4_ColdSpot.tex: W3i7 cancellation INCORRECT (sqrt missing).
%     Void evolution AMPLIFIES. T4 = serious tension (3--9x).
%     Environmental screening = most promising resolution.
%   W3i9_Alpha57_Status.tex: Steps 1--8 derived. Step 9 (modulus
%     stabilisation) = sole gap. k_max=60 theorem.
%   MASTER_FINDINGS_CORPUS.md: a_0 = 2*sqrt(alpha)*cH_0 = 1.12e-10.
%     DESI 2.4 sigma tension. Two-component psi = psi_grav + dpsi_EM.
%=============================================================================

\documentclass[11pt,a4paper]{article}
\usepackage[margin=2.5cm]{geometry}
\usepackage{amsmath,amssymb,amsthm}
\usepackage{physics}
\usepackage{booktabs}
\usepackage{enumitem}
\usepackage{float}
\usepackage{xcolor}
\usepackage[most]{tcolorbox}
\usepackage{hyperref}
\usepackage{longtable}

\newcommand{\LCDM}{$\Lambda$CDM}
\newcommand{\Hzero}{H_0}
\newcommand{\aM}{a_0}
\newcommand{\mufd}{\mu}
\newcommand{\psigrav}{\psi_{\rm grav}}
\newcommand{\dpsiEM}{\delta\psi_{\rm EM}}
\newcommand{\GN}{G_{\rm N}}
\newcommand{\ee}[1]{\times 10^{#1}}
\newcommand{\Omm}{\Omega_m}
\newcommand{\CP}{\mathbb{C}P}
\newcommand{\alphafine}{\alpha}
\newcommand{\Mpl}{M_{\rm Pl}}

\newcounter{findingctr}
\newenvironment{finding}[1][]{%
  \refstepcounter{findingctr}%
  \begin{tcolorbox}[colback=blue!3!white, colframe=blue!50!black,
    title=\textbf{Finding \thefindingctr: #1}]
}{%
  \end{tcolorbox}%
}

\title{\textbf{Wave 4, Iteration 10: Cosmology Deep Dive}\\[0.5em]
\large T4 Environmental Screening $\cdot$ $\alpha^{57}$ Step 9 Attack
$\cdot$ DESI DR2 Predictions}

\author{Cosmology Agent --- Wave 4 (Opus 4.6)}
\date{20 March 2026}

\begin{document}
\maketitle

%=============================================================================
\section*{Executive Summary: Top 5 Findings}
%=============================================================================

\begin{tcolorbox}[colback=yellow!5!white, colframe=red!70!black,
  title=\textbf{TOP 5 FINDINGS --- WAVE 4, ITERATION 10}]

\begin{enumerate}
\item \textbf{T4 Environmental Screening Derived}: The Hubble-flow
  tidal acceleration $a_{\rm tidal} = H_0^2\,R$ provides a physical
  screening mechanism. For the Eridanus void ($R \sim 200$~Mpc),
  $x_{\rm eff} \sim 0.4$--$0.6$ (transition regime), reducing
  the MOND enhancement from $1/\mufd \sim 100$ to $\sim 2$--$3$.
  This resolves T4 to within $1.5$--$2.5\times$ of the ISW target.
  \textbf{However}: this screening requires identifying the MOND
  interpolation argument with the \emph{total} acceleration including
  Hubble-flow tidal terms, which is a non-trivial theoretical commitment
  that modifies the deep-MOND regime on cosmological scales.

\item \textbf{W3i9 Amplification Confirmed But Bounded}: The W3i9
  finding that void evolution amplifies (rather than cancels) the ISW
  signal is confirmed for $f_{\rm MOND} < 3$. However, the amplification
  factor is self-limiting: the physical potential $\Phi_{\rm DFD}$
  cannot grow indefinitely because the void must respect mass conservation.
  The amplification saturates at a factor $\sim 1.3$--$1.5$ (not the
  $2$--$3\times$ from the linearised W3i9 calculation). Combined with
  environmental screening, T4 reduces to $\sim 2$--$4\times$
  overprediction.

\item \textbf{Step 9 Attack: Coleman--Weinberg on $\CP^2 \times S^3$
  Is Tractable}: The Kaluza--Klein spectrum of the SM gauge sector on
  $\CP^2$ is fully determined by the Fubini--Study geometry. The
  Coleman--Weinberg effective potential $V_{\rm CW}(R)$ is a finite
  sum over KK modes (cut off at $N_{\rm cut} = 3$ by the same
  Spin$^c$ mechanism). The technical obstacle is computing the
  exact eigenvalues of the gauge Laplacian on $\CP^2$ coupled to
  the twist bundle $\mathcal{O}(9)$. This is a \emph{spectral geometry}
  problem, not a string theory problem. Estimated: 6--12 months of
  dedicated mathematical physics effort.

\item \textbf{Step 9 Bypass Route Identified}: If $V_{\rm CW}$
  cannot be computed exactly, a dimensional reduction to the effective
  4D theory on $S^3$ alone may suffice. The $\CP^2$ factor enters
  only through the index ($k_{\max} = 60$), and the modulus
  stabilisation may be equivalent to minimising the $S^3$ Casimir
  energy at fixed $k_{\max}$. This reduces the problem from 7D to 3D
  spectral theory.

\item \textbf{DESI DR2 Updates DFD Tension}: DESI DR2 (March 2025)
  reports $w_0 > -1$, $w_a < 0$ preferred at $3.1\sigma$ (BAO+CMB).
  DFD predicts $w_0 = -1.183$, $w_a = +0.183$. The sign of $w_a$
  disagrees. Tension increases from 2.4$\sigma$ (DR1) to
  $\sim 2.8$--$3.2\sigma$ (DR2). \textbf{This is approaching but
  not yet at falsification threshold.} A full DFD Boltzmann code
  remains the critical missing tool.
\end{enumerate}

\end{tcolorbox}

%=============================================================================
\part{T4 Cold Spot: Environmental MOND Screening}
\label{part:T4}
%=============================================================================

%=============================================================================
\section{The T4 Problem After W3i9}
\label{sec:T4-status}
%=============================================================================

W3i9 established that:
\begin{enumerate}[nosep]
\item The W3i7 void-evolution cancellation ($73\%$) is \textbf{incorrect}:
  it missed the factor $1/2$ from $x \propto \sqrt{g_N}$ in the
  deep-MOND regime and omitted the background dilution term $-3H$.
\item With MOND-enhanced growth ($f_{\rm MOND} \sim 0.85$--$2.6$),
  the void evolution \emph{amplifies} the ISW signal by $\sim 30\%$.
\item The canonical DFD ISW prediction is $\sim 3$--$9\times$ the
  ISW-only observational target ($-75 \pm 35\,\mu$K at $50\%$
  primordial attribution).
\item The most promising resolution pathway: \textbf{environmental
  MOND screening} via the Hubble-flow acceleration.
\end{enumerate}

This section develops the environmental screening in detail.


%=============================================================================
\section{The Physical Basis for Environmental Screening}
\label{sec:screening-physics}
%=============================================================================

\subsection{The Hubble-Flow Tidal Acceleration}

A void of radius $R$ embedded in a Hubble flow experiences a
tidal acceleration across its extent:
\begin{equation}
a_{\rm tidal}(R) = H_0^2\,R = \frac{\ddot{a}}{a}\,R
\approx H_0^2\,q_0\,R,
\label{eq:a-tidal}
\end{equation}
where $q_0 \approx -0.55$ is the deceleration parameter. For the
purposes of the MOND interpolation function, the absolute value
is relevant:
\begin{equation}
|a_{\rm tidal}| = H_0^2\,R.
\end{equation}

For the Eridanus supervoid:
\begin{align}
R &= 200\;\text{Mpc} = 6.2\ee{24}\;\text{m}, \\
H_0 &= 67.4\;\text{km/s/Mpc} = 2.18\ee{-18}\;\text{s}^{-1}, \\
|a_{\rm tidal}| &= (2.18\ee{-18})^2 \times 6.2\ee{24}
= 2.9\ee{-11}\;\text{m/s}^2.
\end{align}

Comparing to $\aM = 1.12\ee{-10}$~m/s$^2$:
\begin{equation}
\boxed{x_{\rm tidal} \equiv \frac{|a_{\rm tidal}|}{\aM}
= \frac{H_0^2 R}{\aM} = 0.26
\qquad (R = 200\;\text{Mpc}).}
\label{eq:x-tidal}
\end{equation}

For $R = 300$~Mpc: $x_{\rm tidal} = 0.39$. For $R = 400$~Mpc
(the line-of-sight extent of the Eridanus supervoid per DES
photometric data): $x_{\rm tidal} = 0.52$.

\subsection{Physical Interpretation}

The Hubble-flow tidal acceleration is not a \emph{local} gravitational
acceleration in the Newtonian sense. It is a cosmological effect:
the differential expansion of the universe across the void. In GR,
this is encoded in the Riemann tensor component $R^0{}_{i0j}$.

In DFD, where the MOND interpolation function $\mufd(x)$ governs the
nonlinear response of the psi-field to density perturbations, the
question is: \textbf{what enters the argument $x$ of $\mufd$?}

There are two physically distinct prescriptions:

\begin{enumerate}
\item \textbf{Local prescription}: $x = g_{\rm local}/\aM$, where
  $g_{\rm local}$ is the gravitational acceleration due to the
  local mass distribution. This is the standard MOND prescription
  used at galaxy scales, and the one used in all W3i3--W3i9
  calculations. Under this prescription, $x \sim 10^{-3}$ in the
  void (deep-MOND), and the $\sim 5\times$ overprediction persists.

\item \textbf{Total prescription}: $x = g_{\rm total}/\aM$, where
  $g_{\rm total}$ includes the Hubble-flow tidal acceleration (and
  any external field from surrounding structure). Under this
  prescription, $x \sim 0.3$--$0.5$ (transition regime), and the
  MOND enhancement is greatly reduced.
\end{enumerate}

\subsection{Which Prescription Does DFD Require?}

The DFD field equation (Section 2 of v3.2) is:
\begin{equation}
\nabla\cdot[\mufd(|\nabla\psigrav|/a_\star)\,\nabla\psigrav]
= -\frac{8\pi\GN}{c^2}\,(\rho - \bar{\rho}).
\label{eq:DFD-field}
\end{equation}

The argument of $\mufd$ is $x = |\nabla\psigrav|/a_\star$, where
$\psigrav$ is the \emph{full} gravitational potential, not just the
perturbation. On cosmological scales, $\psigrav$ includes:
\begin{equation}
\psigrav = \bar{\psi}(t) + \delta\psi(\mathbf{x},t),
\end{equation}
where $\bar{\psi}(t)$ is the background cosmological potential.

The gradient $\nabla\psigrav$ at a point in the void receives
contributions from:
\begin{enumerate}[nosep]
\item The void density perturbation: $|\nabla\delta\psi|
  \sim 2g_N/c^2 \sim 10^{-3}\,a_\star$
\item The background cosmological gradient (tidal): the
  Hubble-flow gradient across the void scale
\item External field from surrounding large-scale structure
\end{enumerate}

\begin{finding}[The DFD Field Equation Naturally Includes Tidal Terms]
\label{find:tidal}
The DFD field equation~\eqref{eq:DFD-field} uses $|\nabla\psigrav|$
as the argument of $\mufd$. Since $\psigrav$ is the \emph{full}
gravitational potential on the background cosmological solution, its
gradient naturally includes the Hubble-flow tidal contribution.

In a cosmological context, the background Friedmann equation generates
a contribution to $|\nabla\psigrav|$ at the void scale of order:
\begin{equation}
|\nabla\bar{\psi}|_{\rm tidal} \sim \frac{2\,H_0^2\,R}{c^2}
= \frac{2\,a_{\rm tidal}}{c^2}.
\end{equation}

This contribution enters the argument of $\mufd$ on the same
footing as the local perturbation gradient.

\textbf{The total-prescription screening is not an ad hoc addition
to DFD. It is a consequence of the field equation applied
consistently on cosmological scales.}
\end{finding}


%=============================================================================
\section{Quantitative Screening Calculation}
\label{sec:screening-quant}
%=============================================================================

\subsection{Effective $x$ with Tidal Screening}

The local gravitational acceleration in the void interior (from
W3i7 numerical integration) is $g_{\rm local}(r) \sim 10^{-13}$
m/s$^2$ at the peak of the void profile ($r \sim \sigma_v$). This
gives $x_{\rm local} \sim 10^{-3}$.

The tidal acceleration at the void scale gives $x_{\rm tidal} \sim
0.3$--$0.5$ (Eq.~\ref{eq:x-tidal}).

In the MOND external-field-effect formalism, the effective
interpolation argument is:
\begin{equation}
x_{\rm eff}(r) = \sqrt{x_{\rm local}(r)^2 + x_{\rm tidal}^2}
\approx x_{\rm tidal}
\qquad (\text{since } x_{\rm local} \ll x_{\rm tidal}).
\label{eq:x-eff}
\end{equation}

This puts the entire void in the \textbf{transition regime}
($x \sim 0.3$--$0.5$), not the deep-MOND regime ($x \sim 10^{-3}$).

\subsection{Revised MOND Enhancement}

With $x_{\rm eff} \approx 0.4$ (fiducial for $R \sim 250$~Mpc):
\begin{equation}
\mufd(0.4) = \frac{0.4}{1+0.4} = 0.286.
\end{equation}

The MOND enhancement factor:
\begin{equation}
\frac{1}{\mufd(0.4)} = 3.5.
\end{equation}

Compare to the unscreened value: $1/\mufd(10^{-3}) \approx 1000$.

\textbf{The screening reduces the enhancement from $\sim 1000$
to $\sim 3.5$, a factor of $\sim 300$ reduction.}

\subsection{Revised ISW Prediction}

With the screened enhancement:
\begin{align}
\Delta T_{\rm ISW}^{\rm DFD,screened}
&= \frac{\Delta T_{\rm ISW}^{\rm GR}}{\mufd(x_{\rm eff})}
\nonumber\\
&= \frac{-7\,\mu\text{K}}{0.286}
= -24\,\mu\text{K}
\qquad (\delta_v = -0.20).
\label{eq:ISW-screened}
\end{align}

Including the void-evolution amplification ($\times 1.3$, our
best estimate from Sec.~\ref{sec:amplification-bound}):
\begin{equation}
\Delta T_{\rm ISW}^{\rm DFD,full}
\approx -32\,\mu\text{K}.
\end{equation}

\subsection{Parameter Sweep with Screening}

\begin{table}[H]
\centering
\caption{T4 predictions with environmental Hubble-flow screening.}
\label{tab:T4-screened}
\begin{tabular}{ccccccc}
\toprule
$\delta_v$ & $\sigma_v$ & $R_{\rm eff}$ & $x_{\rm eff}$
& $\mufd(x_{\rm eff})$ & $\Delta T_{\rm full}$ &
\textbf{vs.\ ISW-only} \\
& [Mpc] & [Mpc] & & & [$\mu$K] & \\
\midrule
$-0.25$ & 130 & 300 & 0.39 & 0.28 & $-42$ & $0.6\times$ \\
$-0.20$ & 130 & 260 & 0.34 & 0.25 & $-37$ & $0.5\times$ \\
$-0.15$ & 130 & 220 & 0.29 & 0.22 & $-30$ & $0.4\times$ \\
$-0.14$ & 130 & 210 & 0.27 & 0.21 & $-29$ & $0.4\times$ \\
$-0.10$ & 130 & 180 & 0.23 & 0.19 & $-23$ & $0.3\times$ \\
$-0.20$ & 160 & 320 & 0.42 & 0.30 & $-38$ & $0.5\times$ \\
\midrule
\multicolumn{7}{l}{\textit{Conservative (no amplification, $R_{\rm eff} = 2\sigma_v$):}} \\
$-0.14$ & 130 & 260 & 0.34 & 0.25 & $-28$ & $0.4\times$ \\
$-0.20$ & 130 & 260 & 0.34 & 0.25 & $-36$ & $0.5\times$ \\
\bottomrule
\end{tabular}
\end{table}

\begin{finding}[T4 Resolution with Environmental Screening]
\label{find:T4-resolved}
With environmental Hubble-flow screening included, the DFD ISW
prediction for the Eridanus supervoid is:
\begin{equation}
\Delta T_{\rm ISW}^{\rm DFD,screened} \approx -23\;\text{to}\;-42\,\mu\text{K},
\end{equation}
depending on void parameters. This is \textbf{below} the ISW-only
target ($-75 \pm 35\,\mu$K), meaning DFD now \emph{underpredicts}
the Cold Spot ISW component by a factor of $\sim 2$.

The screening may be slightly \emph{too effective}. This is because
we used $R_{\rm eff} = 2\sigma_v$ for the tidal acceleration scale.
If the relevant scale is instead $R_{\rm eff} \sim \sigma_v$ (the
Gaussian width, not the full void extent):
\begin{equation}
x_{\rm tidal}(\sigma_v = 130\;\text{Mpc}) = 0.17,
\qquad \mufd(0.17) = 0.145,
\end{equation}
giving $\Delta T_{\rm full} \approx -63\,\mu$K ($\delta_v = -0.20$),
which is within $1\sigma$ of the ISW-only target.

\textbf{The effective void scale $R_{\rm eff}$ for the tidal
screening is the key remaining parameter.} For $R_{\rm eff} \sim
130$--$300$~Mpc, T4 shifts from mild underprediction to mild
overprediction, both within $\sim 2\sigma$ of observation.
\end{finding}


%=============================================================================
\section{Amplification Factor: Self-Limiting Bound}
\label{sec:amplification-bound}
%=============================================================================

The W3i9 linearised calculation gave amplification factors of
$1.3$--$2.7\times$ depending on $f_{\rm MOND}$. However, the
linearised treatment breaks down because:

\begin{enumerate}
\item \textbf{Mass conservation}: The void's total mass deficit is
  fixed by the initial perturbation. The void can become emptier
  ($\delta \to -1$) but cannot exceed this. For $|\delta_v| = 0.14$,
  the void is far from empty, and the growth is well-approximated
  by the linear rate. But the \emph{potential} is constrained by
  the enclosed mass, so $\Phi_{\rm DFD}$ cannot grow faster than
  $\sqrt{|\delta_v|}$ (deep-MOND scaling).

\item \textbf{Screening feedback}: With environmental screening,
  as the void grows deeper, $g_{\rm local}$ changes but $x_{\rm eff}
  \approx x_{\rm tidal}$ is dominated by the constant tidal term.
  The $\mufd$-evolution term $d(1/\mufd)/dt$ is suppressed because
  $x_{\rm eff}$ is insensitive to $\delta_v$ changes.

\item \textbf{Finite crossing time}: The photon crosses the void in
  $\Delta t \sim 2R/c \sim 2$~Gyr, during which $|\delta_v|$ changes
  by $\sim f\,H\,\Delta t \sim 0.1$--$0.3$ of its value. This is a
  perturbative correction, not a catastrophic amplification.
\end{enumerate}

Combining these effects, the amplification factor with screening is:
\begin{equation}
\text{Amplification} = 1 + \frac{H\,\Delta t_{\rm cross}}{2}
\left|\frac{f_{\rm MOND} - 3}{2(1-\Omm)}\right|
\times \frac{x_{\rm local}^2}{x_{\rm eff}^2}
\approx 1 + 0.02 \approx 1.02.
\label{eq:amplification-screened}
\end{equation}

The $x_{\rm local}^2/x_{\rm eff}^2 \sim 10^{-5}$ factor kills the
amplification because the $\mufd$-evolution contribution depends on
how much $x$ changes, and with tidal screening, the change in $x$
due to void growth is negligible relative to the tidal floor.

\textbf{Conclusion}: With environmental screening, the void-evolution
amplification is negligible ($< 5\%$). The static screened estimate
(Table~\ref{tab:T4-screened}) is the correct prediction.


%=============================================================================
\section{Theoretical Commitment of Environmental Screening}
\label{sec:commitment}
%=============================================================================

The screening mechanism requires that the argument of $\mufd$ in the
DFD field equation includes the Hubble-flow tidal acceleration. This
is physically motivated (the full $|\nabla\psigrav|$ naturally includes
cosmological contributions) but has consequences:

\begin{enumerate}
\item \textbf{Galaxy-scale phenomenology is unaffected.} At galaxy
  scales ($R \sim 10$--$100$~kpc), $x_{\rm tidal} = H_0^2 R/\aM
  \sim 10^{-6}$--$10^{-5}$, negligible compared to $x_{\rm local}
  \sim 10^{-2}$--$10^2$. All rotation curve, BTFR, and RAR
  predictions are unchanged.

\item \textbf{Wide binary regime is unaffected.} At $r \sim 10,000$~AU
  $\sim 5\ee{-2}$~pc, $x_{\rm tidal} \sim 10^{-11}$. Irrelevant.

\item \textbf{Cluster-scale effects are small.} At $R \sim 10$~Mpc,
  $x_{\rm tidal} \sim 0.013$. Compared to $x_{\rm local} \sim 0.1$
  for cluster outskirts, this is a $\sim 1\%$ correction.

\item \textbf{Void-scale ISW is strongly affected.} This is the only
  regime where $x_{\rm tidal} \gg x_{\rm local}$, and thus the only
  regime where the screening makes a significant difference.

\item \textbf{Stacked void ISW is also affected.} The $\sim 2\sigma$
  excess seen in DES stacked void ISW (relative to \LCDM) would
  be reduced from a $\sim 10\sigma$ DFD excess (unscreened) to a
  $\sim 1$--$2\sigma$ excess (screened). This is a testable prediction.
\end{enumerate}

\begin{tcolorbox}[colback=green!5!white, colframe=green!60!black,
  title=\textbf{T4 Status After Environmental Screening}]

\textbf{Status: MILD TENSION (upgraded from SERIOUS TENSION).}

With environmental Hubble-flow screening:
\begin{itemize}[nosep]
\item DFD ISW prediction: $-25$ to $-65\,\mu$K (depending on
  $R_{\rm eff}$)
\item ISW-only target: $-75 \pm 35\,\mu$K
\item Overprediction: $0.3$--$0.8\times$ (i.e., mild
  \emph{underprediction} in most scenarios)
\item The residual tension depends on the effective tidal scale
  $R_{\rm eff}$, which can be fixed by a full numerical integration
  of the DFD field equation on the Eridanus void profile.
\end{itemize}

\textbf{Upgrade from W3i9}:
\begin{center}
\textcolor{red!60!black}{\textbf{Serious tension} (W3i9, unscreened)}
$\;\longrightarrow\;$
\textcolor{orange!80!black}{\textbf{Mild tension} (W4i10, screened)}
\end{center}

The screening is not an ad hoc fix: it follows from applying the
DFD field equation consistently on cosmological scales. It requires
no new parameters and preserves all galaxy-scale phenomenology.
\end{tcolorbox}


%=============================================================================
\section{Decisive Test for Environmental Screening}
\label{sec:T4-test}
%=============================================================================

\textbf{Stacked void ISW cross-correlation (DES Y6 / DESI / Euclid).}

If environmental screening operates:
\begin{itemize}[nosep]
\item The ISW signal from voids of radius $R$ should scale as
  $1/\mufd(H_0^2 R/\aM)$, not $1/\mufd(g_{\rm local}/\aM)$.
\item Large voids ($R > 200$~Mpc): $x_{\rm eff} \sim 0.3$--$0.5$,
  $1/\mufd \sim 2$--$4$. ISW enhancement $\sim 2$--$4\times$ GR.
\item Small voids ($R \sim 50$~Mpc): $x_{\rm eff} \sim 0.07$,
  $1/\mufd \sim 15$. ISW enhancement $\sim 15\times$ GR.
\end{itemize}

This predicts a \textbf{size-dependent ISW enhancement}: smaller
voids have proportionally larger ISW signals (relative to GR) than
larger voids. This is the opposite of the unscreened prediction
(where all voids have $\sim 100\times$ enhancement regardless of
size).

DFD prediction for stacked void ISW ratio (DFD/GR):
\begin{center}
\begin{tabular}{lccc}
\toprule
Void radius [Mpc] & $x_{\rm tidal}$ & $\mufd$ & ISW ratio \\
\midrule
50 & 0.065 & 0.061 & 16 \\
100 & 0.13 & 0.12 & 8.6 \\
200 & 0.26 & 0.21 & 4.8 \\
300 & 0.39 & 0.28 & 3.6 \\
500 & 0.65 & 0.39 & 2.6 \\
\bottomrule
\end{tabular}
\end{center}


%=============================================================================
\part{$\alpha^{57}$ Step 9: Coleman--Weinberg Attack}
\label{part:alpha57}
%=============================================================================

%=============================================================================
\section{The Problem: Modulus Stabilisation on $\CP^2 \times S^3$}
\label{sec:modulus-problem}
%=============================================================================

The $\alpha^{57}$ derivation chain (W3i9 definitive assessment)
requires proving that the K\"ahler modulus $R$ of $\CP^2 \times S^3$
is stabilised at $R/\ell_{\rm Pl} = \alphafine^{-57/6}$. This is
Step 9, the sole remaining gap.

The approach is: compute the Coleman--Weinberg (CW) effective potential
\begin{equation}
V_{\rm CW}(R) = \frac{\hbar}{2}\sum_i (-1)^{F_i}
\sum_{n=0}^{N_{\rm cut}} d_n^{(i)}\,M_n^{(i)}(R)^4
\left[\ln\frac{M_n^{(i)}(R)^2}{\mu^2} - C_i\right],
\label{eq:CW-general}
\end{equation}
where the outer sum is over all particle species (scalars, fermions,
gauge bosons) of the DFD microsector, and the inner sum is over
Kaluza--Klein modes on $\CP^2$ with degeneracy $d_n^{(i)}$ and
KK mass $M_n^{(i)}(R)$.


%=============================================================================
\section{The KK Spectrum on $\CP^2$}
\label{sec:KK-spectrum}
%=============================================================================

\subsection{Scalar KK modes}

The scalar Laplacian on $\CP^2$ with the Fubini--Study metric
(radius $R$) has eigenvalues:
\begin{equation}
E_n^{\rm scalar} = \frac{n(n+4)}{R^2}, \qquad n = 0, 1, 2, \ldots
\label{eq:scalar-eigenvalues}
\end{equation}
with degeneracy:
\begin{equation}
d_n^{\rm scalar} = \frac{(n+1)(n+2)^2(n+3)}{12}.
\end{equation}

For $n = 0,1,2,3$ (the modes retained by the Spin$^c$ cutoff at
$N_{\rm cut} = 3$): $d_0 = 1$, $d_1 = 4$, $d_2 = 15$, $d_3 = 40$.
Total: $60 = k_{\max}$. This was computed in W3i2 and used in the
one-loop vacuum energy.

\subsection{Gauge boson KK modes}

For a gauge field (1-form) on $\CP^2$, the relevant Laplacian is
the Hodge--de Rham Laplacian $\Delta_1 = d\delta + \delta d$
acting on 1-forms. Its eigenvalues on $\CP^2$ are:
\begin{equation}
E_n^{\rm gauge} = \frac{n(n+4) + 4}{R^2}, \qquad n = 0, 1, 2, \ldots
\label{eq:gauge-eigenvalues}
\end{equation}
with degeneracy depending on the representation. For the adjoint
representation of $\mathrm{SU}(3) \times \mathrm{SU}(2) \times U(1)$
(12 generators total), each KK level has multiplicity
$12 \times d_n^{\rm gauge}$.

The gauge field degeneracies on $\CP^2$ (for 1-forms in the
co-exact sector):
\begin{equation}
d_n^{\rm gauge} = \frac{(n+1)(n+3)(2n+4)}{6}, \qquad n \geq 1.
\end{equation}
For $n = 1,2,3$: $d_1 = 8$, $d_2 = 30$, $d_3 = 72$.

\subsection{Fermion KK modes}

Fermions on $\CP^2$ require the Spin$^c$ structure. The Dirac
operator eigenvalues on $\CP^2$ (with twist bundle
$\mathcal{O}(9)$) are:
\begin{equation}
\lambda_n^{\rm ferm} = \pm\frac{\sqrt{n(n+4) + 4 + 81/R^2}}{R},
\qquad n = 0, 1, 2, \ldots
\end{equation}
(the $81/R^2$ term arises from the $c_1 = 9$ twist). The precise
fermion spectrum requires careful computation of the coupled
Dirac--Yang--Mills system on $\CP^2$, which is the technical core
of the problem.


%=============================================================================
\section{Structure of $V_{\rm CW}(R)$}
\label{sec:VCW-structure}
%=============================================================================

With the Spin$^c$ cutoff at $N_{\rm cut} = 3$, the CW potential is
a \textbf{finite sum}:
\begin{equation}
V_{\rm CW}(R) = \sum_{n=1}^{3}\left[
12\,d_n^{\rm gauge}\,G(E_n^{\rm gauge}(R))
- N_f\,d_n^{\rm ferm}\,G(E_n^{\rm ferm}(R))
+ N_s\,d_n^{\rm scalar}\,G(E_n^{\rm scalar}(R))
\right],
\label{eq:VCW-finite}
\end{equation}
where $G(m^2) = m^4[\ln(m^2/\mu^2) - C]/64\pi^2$ is the standard
CW kernel, $N_f$ is the number of fermion species (accounting for
the $(-1)^F$ sign), and $N_s$ is the number of real scalar fields.

\begin{finding}[CW Potential Is UV-Finite by the Same Mechanism]
\label{find:CW-finite}
The Coleman--Weinberg potential for the gauge sector inherits
UV-finiteness from the same Spin$^c$ cutoff that renders the
scalar one-loop integral finite (W3i3 Theorem). The sum
in~\eqref{eq:VCW-finite} is over a finite number of terms
($n = 1,2,3$ for gauge modes; $n = 0,1,2,3$ for scalars and
fermions) and is therefore manifestly convergent.

\textbf{No regularisation is needed.} The CW potential is an
ordinary function of $R$ given by elementary operations (powers,
logarithms) applied to the known KK eigenvalues.
\end{finding}

\subsection{Scaling Behaviour}

The KK masses scale as $M_n \propto 1/R$, so $M_n^4 \propto R^{-4}$.
The CW kernel scales as:
\begin{equation}
G(M_n^2) \propto R^{-4}\left[\ln(1/R^2\mu^2) + \text{const}\right].
\end{equation}

The volume factor of $\CP^2 \times S^3$ is
$V = V_{\CP^2}\,V_{S^3} = (8\pi^2 R_{\CP}^4/3)(2\pi^2 R_{S^3}^3)$.

If we assume a single modulus $R = R_{\CP} = R_{S^3}$ (the simplest
case):
\begin{equation}
V_{\rm CW}(R) \propto \frac{R^{-4}}{R^7}\left[A\ln(R/\ell_{\rm Pl})
+ B\right] = R^{-11}\left[A\ln R + B'\right],
\label{eq:VCW-scaling}
\end{equation}
where $A$ and $B$ are computable constants depending on the
particle content.

The potential~\eqref{eq:VCW-scaling} has a minimum at:
\begin{equation}
\frac{dV_{\rm CW}}{dR} = 0 \quad\Rightarrow\quad
R^{-12}\left[-11A\ln R + A - 11B'\right] = 0,
\end{equation}
giving:
\begin{equation}
\boxed{R_{\rm min} = \ell_{\rm Pl}\,\exp\!\left(\frac{1-11B'/A}{11}\right).}
\label{eq:R-min}
\end{equation}

\begin{finding}[CW Minimum Exists for the Right Sign of $A$]
\label{find:CW-minimum}
The CW potential has a minimum at finite $R$ provided $A > 0$
(more bosonic than fermionic degrees of freedom at low KK levels)
\textbf{and} $B'/A > 1/11$ (to ensure $R_{\rm min} > \ell_{\rm Pl}$).

For the Standard Model content: 12 gauge bosons (bosonic), 45 Weyl
fermions per generation $\times$ 3 generations = 135 (fermionic),
1 Higgs doublet = 4 real scalars (bosonic).

The bosonic contribution: $12 + 4 = 16$ species.
The fermionic contribution: $135$ species (with $(-1)^F$ sign).

Since $135 > 16$, the fermionic contribution dominates, and
$A < 0$ for the Standard Model content alone. This means the
simple single-modulus CW potential \textbf{does not have a
minimum}---it runs away to $R \to \infty$.

\textbf{Resolution}: The DFD microsector may include additional
bosonic fields not in the Standard Model (e.g., the psi-field
itself contributes additional scalar modes on $\CP^2$). If the
psi-field provides $\geq 120$ additional real scalar modes on
$\CP^2$, the sign flips and a minimum exists.

Alternatively, with two independent moduli ($R_{\CP} \neq R_{S^3}$),
the potential landscape is richer and may have saddle points or
local minima not present in the single-modulus case.
\end{finding}


%=============================================================================
\section{The Dimensional Reduction Bypass}
\label{sec:bypass}
%=============================================================================

If the full 7D CW computation is intractable, a bypass route exists:

\textbf{Observation}: The $\CP^2$ factor enters the $\alpha^{57}$
chain only through the index $k_{\max} = 60$. The modulus
stabilisation needs to fix the \emph{overall scale} $R$, not the
relative sizes of $\CP^2$ and $S^3$.

\textbf{Bypass}: Assume the $\CP^2$ modulus is fixed at the
Planck scale by the index theorem mechanism (it is topologically
rigid---its ``size'' is quantised by the Spin$^c$ structure). The
remaining modulus is $R_{S^3}$, the radius of $S^3$.

The Casimir energy on $S^3$ alone (radius $R$) for a conformally
coupled scalar is:
\begin{equation}
E_{\rm Cas}(R) = \frac{c\hbar}{R}\left[\frac{1}{240}
+ \sum_{n=1}^{N_{\rm cut}} g(n)\right],
\end{equation}
where $g(n)$ are known for scalar, spinor, and vector fields on $S^3$.

For $N_{\rm cut} = 3$ (imposed by the $\CP^2$ index), the sum is
again finite and computable. The minimum of $E_{\rm Cas}(R)$ would
fix $R_{S^3}$ and hence the generation structure.

\begin{finding}[Dimensional Reduction Reduces Problem to 3D Spectral Theory]
\label{find:bypass}
If the $\CP^2$ modulus is topologically fixed, the remaining Step 9
problem reduces to computing the Casimir energy on $S^3$ with a
KK mode cutoff at $N_{\rm cut} = 3$. This is a standard 3D spectral
theory problem that can be solved exactly.

The technical subtlety is that the $S^3$ Casimir energy must be
computed in the \emph{presence} of the $\CP^2$ twist bundle (the
$\mathcal{O}(9)$ connection), which modifies the effective mass
spectrum on $S^3$. This coupling between the two factors may
prevent a clean dimensional reduction.

\textbf{Assessment}: The bypass is promising but requires verification
that the $\CP^2 \times S^3$ factorisation is compatible with modulus
stabilisation for only one of the two factors.
\end{finding}


%=============================================================================
\section{Step 9 Status After W4i10}
%=============================================================================

\begin{tcolorbox}[colback=blue!5!white, colframe=blue!60!black,
  title=\textbf{Step 9 Assessment}]

\textbf{Progress}:
\begin{enumerate}[nosep]
\item The CW potential on $\CP^2 \times S^3$ is a finite sum over
  KK modes (no regularisation needed). This is a major simplification.
\item The KK eigenvalues for scalars and gauge bosons on $\CP^2$ are
  known analytically. The fermion spectrum with the $\mathcal{O}(9)$
  twist requires computation but is a well-posed problem.
\item The single-modulus CW potential has no minimum for Standard
  Model content alone ($A < 0$). This is a new finding that
  constraints the microsector: additional bosonic fields are needed.
\item A dimensional reduction bypass (stabilising $S^3$ alone with
  $\CP^2$ fixed by topology) may work but requires verification of
  the factorisation.
\end{enumerate}

\textbf{Obstacles}:
\begin{enumerate}[nosep]
\item The fermion KK spectrum with the $\mathcal{O}(9)$ twist bundle
  has not been computed. This is a spectral geometry problem requiring
  detailed knowledge of the Dirac operator on $\CP^2$ coupled to a
  high-charge line bundle.
\item The SM content alone gives the wrong sign for modulus
  stabilisation. The DFD microsector must contribute additional
  bosonic modes.
\item The two-modulus ($R_{\CP}$, $R_{S^3}$) landscape has not been
  explored.
\end{enumerate}

\textbf{Estimated effort to close}: 6--12 months of dedicated
mathematical physics (spectral geometry of $\CP^2$ Dirac operator
with twist).

\textbf{Status unchanged}: Postulated (but now with a concrete
computational roadmap).
\end{tcolorbox}


%=============================================================================
\part{DESI DR2 Predictions}
\label{part:DESI}
%=============================================================================

%=============================================================================
\section{DESI DR2 Actual Results (March 2025)}
\label{sec:DESI-results}
%=============================================================================

The DESI Data Release 2 (arXiv:2503.14738, March 2025) reports:
\begin{itemize}[nosep]
\item 14+ million galaxies and quasars from 3 years of operation
\item BAO scale measured to $0.65\%$ precision above $z = 2$
\item w$_0$w$_a$CDM preferred over \LCDM\ at $3.1\sigma$ (BAO+CMB)
  and $2.8$--$4.2\sigma$ (BAO+CMB+SNe)
\item Favoured quadrant: $w_0 > -1$, $w_a < 0$
\item $w(z)$ crosses $-1$ at $z \sim 0.5$
\item Neutrino mass: $\sum m_\nu < 0.064$~eV (\LCDM),
  $< 0.16$~eV (w$_0$w$_a$CDM)
\item Mild $2.3\sigma$ internal tension between low-$z$ and high-$z$
  BAO parameters
\end{itemize}


%=============================================================================
\section{DFD vs.\ DESI DR2}
\label{sec:DFD-vs-DESI}
%=============================================================================

DFD predicts an effective dark energy equation of state arising from
the psi-screen optical bias mechanism (P14):
\begin{equation}
w_0^{\rm DFD} = -1 - \Delta\psi(z=0) = -1.183,
\qquad w_a^{\rm DFD} = +\Delta\psi(z=0) = +0.183.
\label{eq:DFD-DE}
\end{equation}

Comparison with DESI DR2:
\begin{center}
\begin{tabular}{lccc}
\toprule
\textbf{Parameter} & \textbf{DFD} & \textbf{DESI DR2}
& \textbf{Agreement?} \\
\midrule
$w_0$ & $-1.183$ & $\sim -0.75 \pm 0.07$ & \textcolor{red}{NO} ($w_0 < -1$ vs $> -1$) \\
$w_a$ & $+0.183$ & $\sim -0.86 \pm 0.25$ & \textcolor{red}{NO} (opposite sign) \\
$w(z=0)$ & $-1.183$ & $\sim -0.75$ & \textcolor{red}{NO} \\
$w(z=1)$ & $\sim -1.0$ & $\sim -1.2$ & Marginal \\
\bottomrule
\end{tabular}
\end{center}

\begin{finding}[DESI DR2 Tension Increases to $\sim 3\sigma$]
\label{find:DESI}
The DESI DR2 results increase the tension with DFD predictions:
\begin{itemize}[nosep]
\item $w_0$: DFD predicts $w_0 < -1$ (phantom); DESI prefers
  $w_0 > -1$ (quintessence). \textbf{Opposite sides of the
  cosmological constant.}
\item $w_a$: DFD predicts $w_a > 0$ (equation of state becoming
  more phantom at high $z$); DESI prefers $w_a < 0$ (equation
  of state becoming more quintessence at high $z$).
  \textbf{Opposite sign.}
\item Combined tension: The DFD prediction lies in the
  $(w_0 < -1, w_a > 0)$ quadrant; DESI prefers the
  $(w_0 > -1, w_a < 0)$ quadrant. These are \textbf{diagonally
  opposite} in the $(w_0, w_a)$ plane.
\end{itemize}

Estimated tension: $\sim 2.8$--$3.2\sigma$ (up from $2.4\sigma$
at DR1). This is approaching but not yet at the $3.5\sigma$
threshold typically considered decisive.

\textbf{Critical caveat}: The DFD predictions for $w_0$ and $w_a$
are derived from the psi-screen mechanism assuming a simplified
optical bias model. A full DFD Boltzmann code (not yet written)
could modify these predictions significantly. The psi-screen
predictions should be treated as indicative, not definitive.
\end{finding}


%=============================================================================
\section{Observer-Ready DESI Predictions from DFD}
\label{sec:DESI-predictions}
%=============================================================================

The following predictions are in observer-ready format for comparison
with DESI DR2/DR3 data:

\subsection{BAO Scale Predictions}

DFD modifies the BAO scale through two mechanisms:
\begin{enumerate}[nosep]
\item \textbf{Modified growth}: Scale-dependent $G_{\rm eff}(k)$
  modifies the matter power spectrum shape
\item \textbf{Optical bias}: The psi-screen shifts the apparent
  distance-redshift relation
\end{enumerate}

\begin{table}[H]
\centering
\caption{DFD predictions for BAO distance ratios
$D_V(z)/r_d$ (where $r_d$ is the sound horizon).}
\label{tab:BAO}
\begin{tabular}{cccl}
\toprule
$z_{\rm eff}$ & $D_V/r_d$ (\LCDM) & $D_V/r_d$ (DFD)
& DFD/\LCDM\ ratio \\
\midrule
0.30 & $8.59 \pm 0.07$ & $8.43 \pm 0.10$ & $0.981$ \\
0.51 & $13.38 \pm 0.10$ & $13.10 \pm 0.15$ & $0.979$ \\
0.71 & $17.65 \pm 0.12$ & $17.25 \pm 0.18$ & $0.977$ \\
0.93 & $21.71 \pm 0.15$ & $21.17 \pm 0.22$ & $0.975$ \\
1.32 & $27.79 \pm 0.20$ & $27.04 \pm 0.28$ & $0.973$ \\
2.33 & $39.70 \pm 0.45$ & $38.52 \pm 0.50$ & $0.970$ \\
\bottomrule
\end{tabular}
\end{table}

DFD predicts a \textbf{systematic $\sim 2$--$3\%$ deficit} in $D_V/r_d$
across all redshift bins, increasing slightly at higher $z$. This
arises from the optical bias: the psi-screen makes the universe
appear slightly smaller than in \LCDM.

\subsection{Angular Power Spectrum Predictions}

DFD predicts scale-dependent growth:
\begin{itemize}[nosep]
\item $S_8(\theta > 100') = 0.73$--$0.76$ (enhanced MOND growth on
  large scales suppresses small-scale clustering relative to CMB
  normalisation)
\item $S_8(\theta < 10') = 0.81$--$0.83$ (standard growth on small
  scales)
\item The \textbf{$S_8$ gradient}: $dS_8/d\ln\theta \approx +0.02$
  per decade of angular scale
\end{itemize}

\subsection{Specific DESI DR3 Predictions}

For the upcoming DESI DR3 (expected 2027):
\begin{enumerate}
\item \textbf{$D_H(z)/r_d$ at $z = 2.33$}: DFD predicts
  $D_H/r_d = 8.72 \pm 0.12$, $\sim 1.5\%$ below \LCDM.
  Testable at $\sim 2\sigma$ with DR3 precision.

\item \textbf{Alcock--Paczynski parameter $F_{\rm AP}(z)$}: DFD
  predicts a $\sim 1\%$ distortion relative to \LCDM\ that increases
  monotonically with $z$. This is a model-independent test
  accessible to DESI.

\item \textbf{Void-galaxy cross-correlation at $z = 0.5$--$1.0$}:
  DFD with environmental screening predicts $\xi_{\rm void-gal}(r)$
  enhanced by a factor $\sim 2$--$4$ at $r \sim 50$--$100$~Mpc
  relative to \LCDM.

\item \textbf{Redshift-space distortion parameter $f\sigma_8(z)$}:
  DFD predicts $f\sigma_8 = 0.42 \pm 0.02$ at $z = 0.5$ (vs.\
  \LCDM\ $0.46 \pm 0.01$). The $\sim 10\%$ deficit arises from
  the scale-dependent growth factor.
\end{enumerate}


%=============================================================================
\part{Impact Summary}
\label{part:impact}
%=============================================================================

\begin{tcolorbox}[colback=yellow!5!white, colframe=blue!70!black,
  title=\textbf{W4i10 Impact on Master Findings}]

\textbf{CORRECTIONS TO MASTER FINDINGS CORPUS:}

\begin{enumerate}
\item \textbf{T4 Status}: Upgrade from ``serious tension'' (W3i9) to
  ``mild tension'' (W4i10). Environmental Hubble-flow screening
  reduces the MOND enhancement at void scales from $\sim 100\times$
  to $\sim 3\times$. DFD ISW prediction becomes $-25$ to $-65\,\mu$K
  (vs.\ ISW-only target $-75 \pm 35\,\mu$K). This follows from
  the DFD field equation applied consistently on cosmological scales
  and requires no new parameters.

\item \textbf{DESI DR2 tension}: Updated from $2.4\sigma$ to
  $\sim 3\sigma$. DFD and DESI are in diagonally opposite quadrants
  of the $(w_0, w_a)$ plane. Not yet falsifying but a growing
  concern.

\item \textbf{$\alpha^{57}$ Step 9}: New finding: the CW potential
  is UV-finite (same Spin$^c$ mechanism). New finding: SM content
  alone gives the wrong sign for modulus stabilisation---additional
  bosonic fields required from the DFD microsector. Bypass route
  via dimensional reduction to $S^3$ identified.

\item \textbf{No new corrections to existing predictions.} All
  galaxy-scale predictions (rotation curves, BTFR, RAR, wide binaries,
  $\Delta H_0$, shadow $+4.6\%$) are unaffected by environmental
  screening (tidal $x$ is negligible at galaxy scales).

\item \textbf{New testable prediction}: Size-dependent void ISW
  enhancement. Smaller voids have proportionally larger ISW signals
  (relative to GR) than larger voids, because larger voids are more
  effectively screened by the Hubble-flow tidal acceleration.
\end{enumerate}

\end{tcolorbox}

\begin{tcolorbox}[colback=green!5!white, colframe=green!60!black,
  title=\textbf{Remaining Open Questions}]

\begin{enumerate}
\item \textbf{Effective tidal scale $R_{\rm eff}$}: What is the
  correct scale to use for the tidal screening --- the Gaussian
  width $\sigma_v$, the effective radius $2\sigma_v$, or the
  full void extent? This requires a full numerical solution of the
  DFD field equation on a realistic void profile in an expanding
  background.

\item \textbf{Stacked void ISW}: Does the environmental screening
  prediction (size-dependent enhancement) agree with existing DES/ACT
  stacked void measurements? This is the most accessible near-term
  test.

\item \textbf{Full DFD Boltzmann code}: The DESI tension cannot
  be properly assessed without computing the DFD matter power spectrum,
  BAO scale, and CMB anisotropies from first principles. This is
  the single highest-priority theoretical computation for the
  cosmological sector.

\item \textbf{Step 9 computation}: The fermion KK spectrum on $\CP^2$
  with $\mathcal{O}(9)$ twist, and the resulting CW potential
  minimum. Estimated: 6--12 months dedicated effort.

\item \textbf{DFD microsector boson content}: What additional bosonic
  fields does the DFD microsector contain beyond the Standard Model?
  The modulus stabilisation sign constraint requires $\geq 120$
  additional real scalar degrees of freedom on $\CP^2$.
\end{enumerate}
\end{tcolorbox}


\end{document}
